% ============================================================================
\section{概念：为什么需要系统调用}
\label{sec:syscall-concept}
% ============================================================================

上一章我们成功从 Ring 0 跳转到了 Ring 3。然而用户态代码几乎什么都做不了——
不能输出字符（\texttt{out} 是特权指令）、不能停机（\texttt{hlt} 会 \#GP）、
不能访问内核内存（PTE.US=0 的页会触发 \#PF）。

那用户程序怎么打印 ``Hello World''？答案就是\term{系统调用}（system call）。

\begin{whybox}
系统调用是用户态进入内核的\textbf{唯一受控入口}。
与普通函数调用不同，它涉及：
\begin{enumerate}
  \item \textbf{特权级切换}：CPL 从 3 升到 0，CPU 自动换栈
  \item \textbf{参数验证}：内核必须检查用户传入的指针是否合法
  \item \textbf{安全审计}：内核决定"允许做什么"而非用户
\end{enumerate}
类比：系统调用就像银行柜台——你坐在玻璃窗外面（Ring 3）填单子，
柜员（Ring 0）审核后帮你执行操作。你不能直接进保险库。
\end{whybox}

% ----------------------------------------------------------------------------
\subsection{系统调用 vs 普通函数调用}
% ----------------------------------------------------------------------------

初学者容易把 \texttt{write()} 当成普通的 C 函数调用。
实际上用户态的 \texttt{write()} 是一个\textbf{薄封装}（wrapper），
内部执行的是 \texttt{int \$0x80} 或 \texttt{sysenter} 指令，
触发 CPU 的特权级切换机制。

\begin{center}
\begin{tabular}{lll}
  \toprule
  \textbf{项目} & \textbf{普通函数调用} & \textbf{系统调用} \\
  \midrule
  特权级   & 不变（Ring 3 → Ring 3）  & 切换（Ring 3 → Ring 0） \\
  栈       & 同一个用户栈             & CPU 自动切换到内核栈（TSS.ESP0） \\
  开销     & 几纳秒（\texttt{call}+\texttt{ret}） & 数百纳秒（中断流程/MSR 切换） \\
  参数传递 & 栈或寄存器（cdecl/fastcall） & 寄存器（避免访问不可信用户栈） \\
  安全     & 无检查                   & 内核验证每个参数 \\
  \bottomrule
\end{tabular}
\end{center}

% ----------------------------------------------------------------------------
\subsection{三种系统调用入口机制}
\label{sec:syscall-mechanisms}
% ----------------------------------------------------------------------------

x86 提供了三种从 Ring 3 进入 Ring 0 的机制，它们的目的相同但性能和复杂度不同：

\begin{enumerate}
  \item \texttt{int N}（软中断）——最传统的方式。
    Linux i386 早期使用 \texttt{int 0x80}。
    CPU 查 IDT[0x80]，走完整的中断流程（保存上下文 → 切换栈 → 执行 handler）。
    优点：简单、调试友好；缺点：比较慢（需要查 IDT + 保存完整帧）。

  \item \texttt{sysenter}/\texttt{sysexit}（Intel 快速系统调用）——
    Pentium II 引入。通过 MSR（Model Specific Register）硬编码入口点，
    跳过 IDT 查表和完整帧保存。
    显著（约 10×）快于 \texttt{int}，但编程复杂度更高。

  \item \texttt{syscall}/\texttt{sysret}（AMD64）——
    64 位 long mode 专用。本项目是 32 位 i386，暂不涉及。
\end{enumerate}

\begin{designbox}
我们将 \texttt{int 0x80} 和 \texttt{sysenter} 设计为可切换的\textbf{后端}。
本章先搭框架（接口 + dispatch + table），D-3 实现 \texttt{int 0x80}，
D-4 实现 \texttt{sysenter}。切换只需改 \texttt{KCONFIG\_SYSCALL\_BACKEND}。
\end{designbox}

% ============================================================================
\subsection{系统调用全路径概览}
% ============================================================================

从用户态发起一个 \texttt{write(1, "hello", 5)} 到串口看到输出，经历以下步骤：

\begin{figure}[H]
  \centering
  \begin{tikzpicture}[
    box/.style={draw, rounded corners=3pt, minimum width=3.2cm, minimum height=0.9cm,
                font=\small, align=center},
    userbox/.style={box, fill=blue!8},
    kernbox/.style={box, fill=orange!10},
    asmbox/.style={box, fill=red!8},
    arrow/.style={-{Stealth[length=4pt]}, thick},
    note/.style={font=\scriptsize\itshape, text=gray},
  ]

  % --- Ring 3 区域 ---
  \node[userbox] (app) at (0, 0) {用户程序\\{\footnotesize\texttt{write(1, buf, n)}}};
  \node[userbox] (wrapper) at (0, -1.5) {syscall wrapper\\{\footnotesize\texttt{EAX=4, EBX=1, ECX=buf}}};

  % --- 特权级切换 ---
  \node[asmbox] (stub) at (0, -3.5) {汇编 stub\\{\footnotesize\texttt{int \$0x80} 或 \texttt{sysenter}}};

  % --- Ring 0 区域 ---
  \node[kernbox] (dispatch) at (0, -5.5) {\texttt{syscall\_dispatch()}\\{\footnotesize 查表 \texttt{syscall\_table[nr]}}};
  \node[kernbox] (handler) at (0, -7.2) {\texttt{sys\_write()}\\{\footnotesize 执行内核服务}};

  % Arrows
  \draw[arrow] (app) -- (wrapper) node[right, pos=0.5, note] {inline asm 封装};
  \draw[arrow] (wrapper) -- (stub) node[right, pos=0.5, note] {寄存器传参};
  \draw[arrow, red!60!black, very thick] (stub) -- (dispatch)
    node[right, pos=0.5, font=\small\bfseries, text=red!60!black] {特权级切换}
    node[left, pos=0.5, note] {Ring 3 → Ring 0};
  \draw[arrow] (dispatch) -- (handler) node[right, pos=0.5, note] {\texttt{table[4] → sys\_write}};

  % Return path
  \draw[arrow, dashed, blue!60!black] (handler.west) -- ++(-2.0, 0) |- (app.west)
    node[left, pos=0.25, note, text=blue!60!black] {返回值 → EAX};

  % 分隔线
  \draw[thick, dashed, gray] (-4.5, -2.5) -- (4.5, -2.5)
    node[right, font=\small\bfseries] {Ring 3 / Ring 0 边界};

\end{tikzpicture}

  \caption{系统调用全路径：用户态 → 汇编 stub → 特权级切换 → 内核 handler → 返回}
  \label{fig:syscall-overview}
\end{figure}

这条路径中，本章负责中间的"dispatch 层 + handler"部分。
汇编 stub 和特权级切换由后端（\texttt{int 0x80} / \texttt{sysenter}）在下一章实现。
